\chapter{Background}\label{chap:background}

\textit{Large Language Models (LLMs)} and the \textit{transformer} technology they are based on \autocite{NIPS2017_3f5ee243}, have revolutionized the field of artificial intelligence, enabling machines to generate human-like text, answer complex questions, and perform a wide range of natural language processing tasks. As LLMs become increasingly integrated into societal applications, from chatbots and virtual assistants to content generation and decision-making tools, their deployment environments and output consistency have become critical areas of research \autocite{gallegosBiasFairnessLarge2024}.

This chapter provides the necessary background to understand the methods and findings of this thesis. It covers the state-of-the-art in LLM deployment, the challenges of output variance, and the role of deployment environments in shaping model behavior.

\section{Deployments}

The deployment environment of an LLM refers to the infrastructure and interfaces through which the model is accessed. These environments can significantly influence the model's behavior, output consistency, and refusal rates. This definition of 'deployment' is used to avoid repeating all the deployment environments in the following chapters. This is not an official general term but rather a definition specific to this work. Common deployment environments include:

\subsection{Web Applications}

Web applications, such as Yandex's Alice Chat, provide user-friendly interfaces for interacting with LLMs. These platforms are designed for non-technical users and often include features like \textit{RAG (Retrieval Augmented Generation)} to enhance response quality \autocite{lewisRetrievalaugmentedGenerationKnowledgeintensive2020}. However, the integration of external data sources can introduce variability in outputs, as the model may draw from different contexts depending on the specific deployment \autocite{saad-falconARESAutomatedEvaluation2024}.

\subsection{Cloud APIs}

Cloud-based APIs, like \textit{Yandex Cloud} \footnote{\url{https://aistudio.yandex.ru/en}} or \textit{Scaleway} \footnote{\url{https://www.scaleway.com/en/generative-apis/}}, allow developers to integrate LLMs into their applications. These APIs provide programmatic access to the model, enabling automated and scalable interactions. However, API-based deployments may include additional layers of filtering or moderation, which can affect refusal rates and output consistency \autocite{lipphardtDualStandardsExamining2026}.

\subsection{Local Deployments}

Local deployments involve running LLMs on personal or organizational hardware. This approach offers greater control over the model's environment and reduces dependency on external services. However, local deployments require significant computational resources and may lack the real-time updates and enhancements available in cloud-based environments \autocite{sandriniCloudAssessingBenefits2025}.

\section{Output Variance in LLMs}

Output variance refers to the differences in responses generated by an LLM for the same input prompt. This variance can arise from several factors \autocite{cuiSamePromptDifferent2026,lewisRetrievalaugmentedGenerationKnowledgeintensive2020}:

\begin{itemize}
	\item \textbf{Deployment Environments}: Differences in deployment environments, such as the presence of filtering mechanisms or RAG systems, can influence the model's behavior and output consistency.
	\item \textbf{Input Language and Context}: The language and context of the input prompt can affect the model's response. For example, prompts in different languages or with varying levels of detail may yield different outputs.
\end{itemize}

\section{Measuring Output Variance}

Output variance in this work will be quantified using these metrics/methods:

\begin{itemize}
	\item \textbf{Refusal Rate Analysis}: Tracking the frequency with which a model refuses to answer a prompt can provide insights into its behavior and potential biases \autocite{garcia-ferreroRefusalSteeringFinegrained2025,panPoliticalCensorshipLarge2026}.
	\item \textbf{Cosine Similarity}: By embedding model outputs into vector spaces, researchers can compute cosine similarity scores to measure the similarity between responses \cite{sitikhuComparisonSemanticSimilarity2019}.
	\item \textbf{Language Detection}: Analyzing the language of the model's output can reveal patterns in how the model adapts to different input languages.
\end{itemize}

\section{Related Work}

\subsection{IssueBench Dataset}

The \textit{IssueBench} dataset, introduced by \citeauthor{rottgerIssueBenchMillionsRealistic2026a}, is a comprehensive collection of prompts designed to evaluate the bias and refusal rates of LLMs. It includes prompts across various topics and polarities, making it a valuable resource for studying output variance \cite{rottgerIssueBenchMillionsRealistic2026a}.

\subsection{Topic Modeling and Clustering}

Topic modeling techniques, such as \textit{Latent Dirichlet Allocation (LDA)} or clustering algorithms like \textit{K-Means}, are commonly used to categorize and analyze large sets of text data \cite{grootendorstBERTopicNeuralTopic2022a}. These methods help identify common themes and patterns in LLM outputs, facilitating the study of output variance.

\subsection{Embedding Models}

Embedding models, such as \textit{jina-embeddings-v3} \cite{sturuaJinaembeddingsv3MultilingualEmbeddings2024} or \textit{octen-embedding-8b}, convert text into high-dimensional vectors. These embeddings enable researchers to compute similarity scores and analyze the semantic relationships between different outputs \cite{keraghelWordsComparativeAnalysis2024}.

Understanding the factors that contribute to output variance in LLMs is essential for improving their reliability and consistency. This chapter provided an overview of LLMs, their deployment environments, and the methods used to measure and analyze output variance. The following chapters build on this background to explore the specific findings of this thesis, focusing on the behavior of \textit{Yandex}'s LLMs across different deployments.